
\begin{abstract}
The Vera C. Rubin Observatory is not only an observatory, it is also a distributed production engine for science. During operations, it will move thousands of images, metadata, and alerts through the telescope, camera, network, storage, and processing services every night. Rubin will produce up to 20 TB of raw data per night (10TB on the wire with compression), with an alert stream that reaches millions of events and a target latency of nearly 1 minute \cite{2019ApJ...873..111I}.

In this environment, simple monitoring is not enough. We need to ask why something changed, where latency started, which component is under pressure, and whether the behavior is connected with the observing conditions or the state of the facility. This paper describes the observability journey of the Vera C. Rubin Observatory DevOps team, a multi-year effort to build, refine, and scale a modern observability stack for one of the most ambitious astronomical facilities ever constructed.

The platform is based on open-source components, and the result is, not only a software stack, it is an operations model. Dashboards and alerts are treated as code; labels are part of the interface between teams; and the observability system is continuously improved as commissioning and operations reveal new questions. The paper highlights dashboard creation strategies, discusses the hard parts: configuration complexity, label discipline, storage and retention trade-offs, human adoption, and the need to keep improving the system rather than calling it finished \cite{lsst-it-k8s-cookbook}.
\end{abstract}

